\begin{problem}[The arithmetic local ring $\ZZ_{(p)}$]
Fix a prime number $p$.  Describe $\ZZ_{(p)}$ explicitly, determine its units and unique maximal ideal, compute its residue field, and determine its prime spectrum.
\end{problem}

\paragraph{Why this problem matters.}
This is the simplest arithmetic local scheme and a model for local arithmetic geometry.

\paragraph{Useful ideas before starting.}
The primes of $\ZZ$ contained in $(p)$ are only $(0)$ and $(p)$.

\paragraph{Guided hints.}
A denominator is allowed precisely when it is not divisible by $p$.

\begin{solution}
We have
\[
\ZZ_{(p)}=\{a/b\in\QQ:p\nmid b\}.
\]
A fraction is a unit iff $p\nmid a$, so the unique maximal ideal is $p\ZZ_{(p)}$.  The quotient is $\FF_p$.  The prime-localization correspondence shows
\[
\Spec\ZZ_{(p)}=\{(0),p\ZZ_{(p)}\},
\]
with $(0)$ generic and $p\ZZ_{(p)}$ closed.
\end{solution}

\paragraph{What happened mathematically.}
The global arithmetic line $\Spec\ZZ$ has been zoomed to the generic point and one chosen closed prime.

\paragraph{Extensions.}
Compare $\ZZ_{(p)}$ with the complete local ring $\ZZ_p$.

\paragraph{Provenance.}
Preserves and expands the arithmetic local-ring examples in legacy \emph{Theory of Commutative Algebra 12--13}.
